
\documentclass[sigconf,nonacm,pbalance]{../template/acmart}

\settopmatter{authorsperrow=4}
\setcopyright{none}

\newcommand{\patchpanel}[4]{%
  \begin{minipage}[t]{#1}
    \centering\footnotesize
    #2\par\smallskip
    \includegraphics[width=\linewidth]{figures/#3-frame.png}\par\smallskip
    \includegraphics[height=#4]{figures/#3-patch.png}
  \end{minipage}%
}

\begin{document}

\title{GameSplat}
\subtitle{CS 284 Milestone Report}

\author{Jonathan Cao}
\affiliation{\institution{UC Berkeley}\city{Berkeley}\state{California}\country{USA}}
\author{Laurier Ke}
\affiliation{\institution{UC Berkeley}\city{Berkeley}\state{California}\country{USA}}
\author{Justin Lee}
\affiliation{\institution{UC Berkeley}\city{Berkeley}\state{California}\country{USA}}
\author{Jameson Liu}
\affiliation{\institution{UC Berkeley}\city{Berkeley}\state{California}\country{USA}}

\maketitle

\section*{Change in Scope}

Our original plan was to build an end-to-end pipeline from raw gameplay recordings to Gaussian splats, including automatic UI masking and camera pose recovery.
After proposal feedback, we set aside the raw-footage pipeline and made texture reuse the focus of the project.
We now test whether a well-observed patch can improve reconstruction of a poorly observed copy of the same texture, even when the two patches are on different shapes.
The method is designed to work across games.
For our Minecraft experiments, we use \href{https://github.com/ipechman/MineSplat}{MineSplat}, a mod that captures clean frames and exports the corresponding calibrated camera poses.

\section*{Work Completed}

We train a fixed-pose Splatfacto baseline from the images and camera poses.
We then find repeated patches by matching RootSIFT features across overlapping images and triangulating those features into 3D landmarks.
We align groups of landmarks by rotating, scaling, and translating them, then verify each match by projecting the aligned landmarks into reserved training views.

After verifying the matches, we test whether texture refinement improves the reconstruction.
Our method uses the patch alignment to copy splats from the better-sampled patch into the receiving patch at low opacity.
We keep additions that lower rendering error on the reserved views, then refine the receiving patches.
We compare this with a Gaussian unioning method based on \emph{Splat and Replace} (\href{https://repo-sam.inria.fr/nerphys/splat-and-replace/splat_and_replace.pdf}{Violante et al., 2025}).
Using the same feature matching procedure, we take the union of the aligned splats from both patches, replace each patch's original splats with a copy of that template, and fit the copies jointly.
The pipeline exports PLYs and flythroughs.

\begin{figure}[H]
  \centering
  \patchpanel{.48\linewidth}{Better-sampled patch}{donor}{16mm}\hfill
  \patchpanel{.48\linewidth}{Receiving patch}{receiver}{16mm}
  \caption{The full frames (top) locate the better-sampled and receiving patches.
  The zoomed crops (bottom) show the matched stone-brick texture.}
  \Description{Two columns with a boxed full frame above each zoomed feature patch.}
  \label{fig:matched-features}
\end{figure}

\section*{Preliminary Results}

Both tests started from the original 30,000-step Splatfacto baseline and ran 2,000 refinement steps with splats outside the edited patches frozen.

For splat transfer, the matched control trained the same receiving splats on the same schedule without adding any splats, so we could measure what the transferred splats contribute.
Ordinary whole-scene training gave us a second comparison for the same step budget.
Mean test PSNR on the receiving patches was 35.69~dB with transfer, 34.65~dB at baseline, 35.21~dB with receiver-only training, and 35.30~dB with whole-scene training.
On validation, the corresponding scores were 35.88, 34.99, 35.47, and 35.87~dB.

For Gaussian unioning, the control refined the original splats in each patch independently using the same crops, loss, and schedule to test what sharing a union template contributes.
Mean test PSNR across both patches was 32.40~dB with unioning, 33.71~dB at baseline, and 33.75~dB with independent refinement.
The validation scores were 36.40, 37.45, and 37.18~dB, respectively.
We suspect a shared template discards useful geometry when different shapes share a texture.

Both methods changed full-frame PSNR by less than 0.01~dB from the baseline.
A run with many high-resolution views showed how little room there is to improve patches the baseline already reconstructs well.

\begin{figure}[H]
  \centering
  \patchpanel{.31\linewidth}{Held-out input}{receiver}{17mm}\hfill
  \patchpanel{.31\linewidth}{Baseline render}{baseline}{17mm}\hfill
  \patchpanel{.31\linewidth}{Refined render}{refined}{17mm}
  \caption{The held-out reference, Splatfacto baseline, and texture-refined render share a camera pose.
  The top row boxes the receiving patch, and the bottom row shows the same crop in all three views.}
  \Description{Three columns showing boxed full frames above zoomed views of the same patch.}
  \label{fig:held-out-comparison}
\end{figure}

\section*{Remaining Work}

We'll test ways to share appearance while fitting each patch's geometry separately.
We'll compare them with splat transfer, Gaussian unioning, and ordinary training using the same checkpoint and number of steps.
We'll vary image resolution and view coverage and evaluate on held-out images across captures and seeds.
Shifted transfers and equal-count additions will help isolate what the match contributes.

\end{document}
